%% SCRAP: architecture/03-architecture/heartbeat-system/segfault-analysis
%% SOURCE: docs/working/architecture/03-architecture/heartbeat-system/segfault-analysis.md
%% STATUS: HISTORICAL
%% FITS: dev-guide/ch-heartbeat
%% EDITORIAL: lifted — prose rewritten to press voice

\section{Heartbeat Segfault: Root Cause and Repair}

A race condition between the main VM thread and the heartbeat worker thread
corrupted the \lstinline{RollingWindowOfTruth} structure. Both threads wrote to
shared state without synchronization, producing index corruption, concurrent
writes to the execution-history array, dangling snapshot pointers, and
ultimately a segmentation fault under stress tests exceeding ten thousand
heartbeat cycles. The fault was systematic rather than incidental: any unlocked
access to \lstinline{rolling_window_record_execution()} would reproduce it.

\subsection{The Race}

The main thread records executions from the inner interpreter while the
heartbeat thread services the same window:

\begin{lstlisting}[language=C]
/* main VM thread, inner interpreter */
rolling_window_record_execution(&vm->rolling_window, word_id);

/* heartbeat worker thread */
rolling_window_service(&vm->rolling_window);
\end{lstlisting}

The record routine updated \lstinline{window_pos}, the
\lstinline{execution_history} array, \lstinline{total_executions}, the warm-up
flag, the snapshot-pending flag, and the adaptive-check accumulator---all
without a lock.

\subsection{Failure Modes}

\begin{itemize}
  \item \textbf{Index corruption.} Both threads read the same
        \lstinline{window_pos}, write the same slot (one overwriting the
        other), and each increment the index, losing one execution and
        desynchronizing the position.
  \item \textbf{Snapshot corruption.} A double-buffer copy in progress on one
        thread reads partially-modified history as another thread records,
        yielding a torn snapshot.
  \item \textbf{Dangling pointer.} A snapshot view held on one thread is
        invalidated when the other triggers adaptive shrinking that frees or
        reallocates the snapshot buffers; the subsequent dereference faults.
\end{itemize}

\subsection{Synchronization Requirements}

Six fields require protection: \lstinline{execution_history},
\lstinline{window_pos}, \lstinline{total_executions}, the two snapshot buffers,
the adaptive-check accumulator, and the warm flag.

The chosen remedy is a single global lock---the existing
\lstinline{vm->tuning_lock}---guarding every
\lstinline{rolling_window_record_execution()} and
\lstinline{rolling_window_service()} call. The overhead is roughly one hundred
CPU cycles per word execution, under one percent of system time, and contention
is rare because the heartbeat runs on a separate one-millisecond cadence. A
reader--writer lock was considered and deferred; the workload is write-heavy
and does not justify the added complexity.

\subsection{The Fix}

Each call site is wrapped in lock and unlock, the lock is confirmed initialized
in \lstinline{vm_init()}, and the API header is annotated to state that
\lstinline{rolling_window_record_execution()} is not thread-safe and must be
called under \lstinline{vm->tuning_lock}.

\begin{lstlisting}[language=C]
sf_mutex_lock(&vm->tuning_lock);
rolling_window_record_execution(&vm->rolling_window, word_id);
sf_mutex_unlock(&vm->tuning_lock);
\end{lstlisting}

\subsection{Validation}

The fix is validated against five criteria: no segmentation fault after 100{,}000
heartbeat cycles; a clean Valgrind run for memory and data races; deterministic
physics output across runs; all 936 existing tests passing; and zero compiler
warnings under \lstinline{-Wall -Werror}. The stress regime includes an
aggressive 100~$\mu$s heartbeat
(\lstinline{make HEARTBEAT_TICK_NS=100000 test}) and a prolonged
million-word execution under concurrent heartbeat.

\subsection{Regression Prevention}

A code-review rule requires any future \lstinline{rolling_window_*()} call to be
lock-wrapped, supported by thread-safe unit tests, stress tests in CI, and
continuous Valgrind checks.

%% TODO(bob): confirm the line numbers cited in the source (vm.c:1160, 672,
%% 738) still correspond after subsequent refactors before promotion.
